\documentclass[10pt]{article}

\usepackage{amsmath,amssymb,amsthm}
\usepackage[a4paper,landscape,margin=0.34in,includehead]{geometry}
\usepackage{array}
\usepackage{booktabs}
\usepackage{enumitem}
\usepackage{microtype}
\usepackage{multicol}
\usepackage{titlesec}
\usepackage{fancyhdr}
\usepackage{draftwatermark}
\usepackage{atbegshi}
\usepackage[hidelinks]{hyperref}

% =========================
% 2-page hard cap
% =========================
\AtBeginShipout{%
  \ifnum\value{page}>2
    \AtBeginShipoutDiscard
  \fi
}

% =========================
% Watermark (optional)
% =========================
\SetWatermarkText{\href{https://qrsntu.org}{QRS@NTU \,|\, qrsntu.org}}
\SetWatermarkScale{0.9}
\SetWatermarkLightness{0.995}

% =========================
% Header
% =========================
\setlength{\headheight}{1pt}
\pagestyle{fancy}
\fancyhf{}
\fancyhead[L]{\normalsize \textbf{MH1201 Linear Algebra II - Cheat Sheet (Full Course)}}
\fancyhead[R]{\normalsize\textbf{\href{https://qrsntu.org}{QRS@NTU}}}
\renewcommand{\headrulewidth}{0pt}
\renewcommand{\footrulewidth}{0pt}
\fancyfoot[C]{\tiny\thepage}

% =========================
% Tight typography
% =========================
\setlength{\parindent}{0pt}
\setlength{\parskip}{0pt}
\setlist{nosep,leftmargin=1.05em}
\setlength{\abovedisplayskip}{1.5pt}
\setlength{\belowdisplayskip}{1.5pt}
\setlength{\abovedisplayshortskip}{1.5pt}
\setlength{\belowdisplayshortskip}{1.5pt}
\setlength{\tabcolsep}{2pt}
\renewcommand{\arraystretch}{0.94}
\setlength{\columnsep}{10pt}

% =========================
% Headings readability
% =========================
\titleformat{\section}
  {\bfseries\normalsize}
  {}
  {0pt}
  {\underline}
\titleformat{\subsection}
  {\bfseries\small}
  {}
  {0pt}
  {}
\titlespacing*{\section}{0pt}{2.5pt}{0.5pt}
\titlespacing*{\subsection}{0pt}{1.5pt}{0.25pt}

% =========================
% Quick macros
% =========================
\newcommand{\R}{\mathbb{R}}
\newcommand{\C}{\mathbb{C}}
\newcommand{\F}{\mathbb{F}}
\newcommand{\N}{\mathbb{N}}
\newcommand{\rank}{\operatorname{rank}}
\newcommand{\nullity}{\operatorname{nullity}}
\newcommand{\Ker}{\operatorname{Ker}}
\newcommand{\Range}{\operatorname{Range}}
\newcommand{\im}{\operatorname{im}}
\newcommand{\spann}{\operatorname{span}}
\newcommand{\tr}{\operatorname{tr}}
\newcommand{\diag}{\operatorname{diag}}
\newcommand{\proj}{\operatorname{proj}}
\newcommand{\algmult}{\operatorname{algmult}}
\newcommand{\geomult}{\operatorname{geomult}}
\newcommand{\ip}[2]{\langle #1,#2\rangle}
\newcommand{\ds}{\displaystyle}

\begin{document}

\begin{multicols}{3}
\scriptsize
\raggedcolumns

%==================================================
\section*{Vector Spaces (Ch1)}

\subsection*{Subspaces / span / sums}
$W\le V\iff W\neq\varnothing$ and $cu+v\in W\ \forall u,v\in W,\ c\in\F$. In particular: $0\in W$, closed under $+$ and scalar mult.\
Fast disproof: $0\notin W\Rightarrow W\not\le V$. Kernel trick: $T$ linear $\Rightarrow \Ker(T)\le V$.\
\[
\spann(S)=\{\text{finite linear combinations of vectors of }S\},\qquad \spann(\varnothing)=\{0\}.
\]
Smallest-subspace property: $S\subseteq W\le V\Rightarrow \spann(S)\subseteq W$.\
For $A,B\subseteq V$:
\[
A\subseteq B\Rightarrow \spann(A)\subseteq\spann(B),\qquad
\spann(A\cup B)=\spann(A)+\spann(B).
\]
If $U,W\le V$, then $U\cap W\le V$ and $U+W=\spann(U\cup W)\le V$.

\subsection*{Direct sums}
\[
V=V_1\oplus\cdots\oplus V_n
\iff
\sum_i v_i=0\ (v_i\in V_i)\Rightarrow v_i=0\ \forall i.
\]
Two-subspace case:
\[
V=U\oplus W\iff V=U+W\text{ and }U\cap W=\{0\}.
\]
Pairwise trivial intersections are \emph{not} enough for $n\ge 3$.\
If $S$ is LI and $S=A\sqcup B$, then $\spann(S)=\spann(A)\oplus\spann(B)$.

%==================================================
\section*{Finite-Dimensional Spaces (Ch2)}

\subsection*{LI / dependence / bases}
$S$ linearly independent (LI) iff
\[
\lambda_1v_1+\cdots+\lambda_kv_k=0\Rightarrow \lambda_1=\cdots=\lambda_k=0.
\]
Dependence criterion:
\[
S\text{ dependent}\iff \exists v\in S:\ v\in\spann(S\setminus\{v\}).
\]
Facts: $\varnothing$ is LI; $\{v\}$ LI $\iff v\neq 0$; LI $\Rightarrow$ all subsets LI; dependent $\Rightarrow$ all supersets dependent.\
For $A=[v_1\ \cdots\ v_n]$ with $v_i\in\F^m$,
\[
\{v_i\}\text{ LI}\iff A\lambda=0\text{ only trivial}\iff \operatorname{rref}(A)\text{ pivot in every col}.
\]
A basis = LI $+$ spanning. If $\beta=(b_1,\dots,b_n)$ is ordered and $v=\sum\lambda_i b_i$, then
\[
[v]_\beta=(\lambda_1,\dots,\lambda_n)^{\mathsf T}.
\]
Order matters.

\subsection*{Dimension logic}
If $\dim V=n$:
\[
|S|=n\Rightarrow S\text{ basis}\iff S\text{ LI}\iff S\text{ spans}.
\]
Also: any LI set has size $\le n$; any spanning set has size $\ge n$.\
Basis extension: LI set extends to a basis. Basis extraction: spanning set contains a basis.\
Add-one-more:
\[
S\text{ LI}\Rightarrow S\cup\{v\}\text{ LI}\iff v\notin\spann(S).
\]
Dimension formulas:
\[
\dim(U+W)=\dim U+\dim W-\dim(U\cap W),
\]
\[
V=U\oplus W\iff V=U+W\text{ and }\dim V=\dim U+\dim W.
\]
Standard dims:
\[
\dim\F^n=n,\quad \dim\F^{m\times n}=mn,\quad \dim P_n(\F)=n+1.
\]
Useful matrix subspaces (over indicated fields):
\[
\dim\{A\in\R^{n\times n}:A^{\mathsf T}=A\}=\tfrac{n(n+1)}2,
\]
\[
\dim\{A\in\R^{n\times n}:A^{\mathsf T}=-A\}=\tfrac{n(n-1)}2,
\]
\[
\dim\{A\in\F^{n\times n}:\tr(A)=0\}=n^2-1.
\]

%==================================================
\section*{Linear Transformations (Ch3)}

\subsection*{Kernel / range / rank-nullity}
$T:V\to W$ linear iff $T(u+v)=T(u)+T(v)$ and $T(cu)=cT(u)$, equivalently
\[
T(\alpha u+\beta v)=\alpha T(u)+\beta T(v).
\]
Then $T(0)=0$, $T(-v)=-T(v)$, and $T$ is determined uniquely by its values on a basis.\
\[
\Ker(T)=\{v:T(v)=0\}\le V,\qquad \Range(T)=\{T(v):v\in V\}\le W.
\]
If $\dim V<\infty$,
\[
\dim V=\rank(T)+\nullity(T),\qquad \rank(T)=\dim\Range(T),\ \nullity(T)=\dim\Ker(T).
\]
Injective tests:
\[
T\text{ inj}\iff \Ker(T)=\{0\}\iff \nullity(T)=0\iff \rank(T)=\dim V.
\]
Surjective:
\[
T\text{ onto }W\iff \Range(T)=W\iff \rank(T)=\dim W.
\]
If $\dim V=\dim W$, then
\[
\text{injective}\iff \text{surjective}\iff \text{bijective}\iff \text{invertible}.
\]
Composition: $(T\circ S)^{-1}=S^{-1}\circ T^{-1}$ when both invertible.

\subsection*{Useful templates}
If $U\le W$, then $T^{-1}(U)=\{v:T(v)\in U\}\le V$.\
Idempotent $P^2=P$:
\[
V=\Ker(P)\oplus\Range(P),\qquad v=(v-Pv)+Pv.
\]
If $\ell:V\to\F$ nonzero linear and $\dim V<\infty$,
\[
\dim\Ker(\ell)=\dim V-1,
\]
so hyperplanes are kernels of nonzero functionals.

%==================================================
\section*{Matrices / Coordinates (Ch4)}

\subsection*{Matrix representations}
Let $\beta=(b_1,\dots,b_n)$ basis of $V$, $\gamma=(c_1,\dots,c_m)$ basis of $W$.
\[
[T]^\gamma_\beta=\bigl[\ [T(b_1)]_\gamma\ \cdots\ [T(b_n)]_\gamma\ \bigr]\in\F^{m\times n},
\]
\[
[T(v)]_\gamma=[T]^\gamma_\beta[v]_\beta.
\]
Composition:
\[
[T\circ S]^\gamma_\alpha=[T]^\gamma_\beta[S]^\beta_\alpha.
\]
If $T$ invertible,
\[
[T^{-1}]^\beta_\gamma=([T]^\gamma_\beta)^{-1}.
\]
RREF criteria for $A=[T]^\gamma_\beta\in\F^{m\times n}$:
\[
T\text{ inj}\iff \operatorname{rref}(A)\text{ pivot in every col},
\]
\[
T\text{ onto }W\iff \operatorname{rref}(A)\text{ pivot in every row}.
\]

\subsection*{Change of basis / similarity}
Coordinate change:
\[
[v]_\gamma=[I_V]^\gamma_\beta[v]_\beta,\qquad [I_V]^\beta_\gamma=([I_V]^\gamma_\beta)^{-1}.
\]
With $P_\beta=[I_V]^\sigma_\beta=\bigl[[b_1]_\sigma\ \cdots\ [b_n]_\sigma\bigr]$,
\[
[I_V]^\gamma_\beta=P_\gamma^{-1}P_\beta.
\]
Same operator, different bases:
\[
[T]^\gamma_\gamma=[I_V]^\gamma_\beta\,[T]^\beta_\beta\,[I_V]^\beta_\gamma.
\]
Similarity:
\[
A\sim B\iff \exists Q\in GL_n(\F):\ B=Q^{-1}AQ.
\]
Invariants under similarity:
\[
\det(B)=\det(A),\qquad \tr(B)=\tr(A),\qquad p_B(\lambda)=p_A(\lambda).
\]
Useful identities:
\[
\det(Q^{-1}AQ)=\det(A),\qquad \tr(XY)=\tr(YX).
\]
For $A=\begin{bmatrix}a&b\\ c&d\end{bmatrix}$,
\[
p_A(\lambda)=\lambda^2-\tr(A)\lambda+\det(A).
\]

%==================================================
\section*{Eigenvalues / Diagonalization (Ch5)}

\subsection*{Eigenvalues and eigenspaces}
\[
\lambda\text{ eigenvalue of }T\iff \Ker(T-\lambda I)\neq\{0\},\qquad E_{T,\lambda}=\Ker(T-\lambda I).
\]
For $A=[T]^\beta_\beta$ on $n$-dim space,
\[
\lambda\text{ eigenvalue}\iff \det(A-\lambda I_n)=0,
\]
\[
p_T(\lambda)=\det(A-\lambda I_n)\quad\text{(basis-independent).}
\]
Distinct eigenvalues $\Rightarrow$ corresponding eigenvectors are LI, and
\[
E_{\lambda_1}+\cdots+E_{\lambda_k}=E_{\lambda_1}\oplus\cdots\oplus E_{\lambda_k}.
\]
Triangular matrix: eigenvalues are diagonal entries (with algebraic multiplicities).\
Field warning: roots must lie in the base field.

\subsection*{Diagonalizability}
$T$ diagonalizable iff $V$ has a basis of eigenvectors iff $[T]^\beta_\beta$ is diagonal for some basis $\beta$.\
If $A=PDP^{-1}$ with columns of $P$ eigenvectors and diagonal of $D$ matching eigenvalues in the same order, then
\[
A^k=PD^kP^{-1},\qquad D^k=\diag(\lambda_1^k,\dots,\lambda_n^k).
\]
Criteria:
\[
T\text{ diagonalizable}\iff \dim V=\sum_\lambda \dim E_{T,\lambda}.
\]
If $p_T$ splits, then
\[
T\text{ diagonalizable}\iff \geomult(\lambda)=\algmult(\lambda)\ \forall\lambda.
\]
Always $1\le \geomult(\lambda)\le \algmult(\lambda)$ for each eigenvalue.\
Sufficient: $n$ distinct eigenvalues in an $n$-dimensional space $\Rightarrow$ diagonalizable.\
Not converses: repeated eigenvalue does \emph{not} imply defective; distinct eigenvalues are sufficient, not necessary.\
Nilpotent shortcut:
\[
T^m=0\text{ for some }m\Rightarrow \text{only eigenvalue is }0.
\]
Hence nonzero nilpotent $\Rightarrow$ not diagonalizable. More generally, if $p(T)=0$ and $p$ splits into \emph{distinct} linear factors, then $T$ is diagonalizable.

\subsection*{Polynomial calculus}
If $Av=\lambda v$, then
\[
p(A)v=p(\lambda)v.
\]
So if $A=PDP^{-1}$ is diagonalizable, then
\[
p(A)=P\,\diag(p(\lambda_1),\dots,p(\lambda_n))\,P^{-1}.
\]

%==================================================
\section*{Inner Product Spaces (Ch6)}

\subsection*{Inner products / orthogonality}
On $V$ over $\R$ or $\C$, $\ip{\cdot}{\cdot}$ must satisfy: conjugate symmetry, sesquilinearity (real case: bilinearity), and positive definiteness.\
Norm:
\[
\|u\|=\sqrt{\ip{u}{u}},\qquad d(u,v)=\|u-v\|.
\]
Cauchy--Schwarz and triangle:
\[
|\ip{u}{v}|\le \|u\|\,\|v\|,\qquad \|u+v\|\le \|u\|+\|v\|.
\]
Equality in Cauchy--Schwarz iff $u,v$ are linearly dependent.\
$u\perp v\iff \ip{u}{v}=0$. Orthogonal nonzero set $\Rightarrow$ LI.\
For orthonormal basis (ONB) $\beta$,
\[
v=\sum_{u\in\beta}\ip{u}{v}u.
\]
For an orthogonal basis (not normalized), coefficients are
\[
v=\sum_i \frac{\ip{u_i}{v}}{\ip{u_i}{u_i}}u_i.
\]
Weighted / matrix inner products commonly used:
\[
\ip{x}{y}=x^{\mathsf T}Gy\ \text{or}\ x^\dagger Hy\quad (G,H\text{ symmetric/Hermitian pos. def.}).
\]
Frobenius inner product:
\[
\ip{A}{B}_F=\tr(A^{\mathsf T}B)=\sum_{i,j}a_{ij}b_{ij}.
\]

\subsection*{Projections / Gram--Schmidt / complements}
Projection onto nonzero $u$:
\[
\proj_u(v)=\frac{\ip{u}{v}}{\ip{u}{u}}u.
\]
Gram--Schmidt from basis $(b_1,\dots,b_n)$:
\[
v_1=b_1,\qquad v_k=b_k-\sum_{j=1}^{k-1}\proj_{v_j}(b_k)
\]
(then normalize $g_k=v_k/\|v_k\|$). Nested spans preserved:
\[
\spann(v_1,\dots,v_k)=\spann(b_1,\dots,b_k).
\]
Orthogonal complement:
\[
W^\perp=\{v\in V:\ip{w}{v}=0\ \forall w\in W\}\le V.
\]
Finite-dimensional facts:
\[
V=W\oplus W^\perp,\qquad \dim W+\dim W^\perp=\dim V,
\]
\[
W\cap W^\perp=\{0\},\qquad (W^\perp)^\perp=W.
\]
If $(b_1,\dots,b_m)$ is an ONB of $W$, then orthogonal projection onto $W$ is
\[
P_W(v)=\sum_{i=1}^m \ip{b_i}{v}b_i,
\]
and $v-P_W(v)\in W^\perp$. This is the unique best approximation to $v$ in $W$.\
Least squares for $Ax\approx y$: if $W=\Range(A)$, then $\hat y=Ab$ is closest to $y$ iff
\[
y-Ab\in W^\perp\iff A^{\mathsf T}(y-Ab)=0\iff A^{\mathsf T}Ab=A^{\mathsf T}y.
\]

\subsection*{Common traps}
Subspace: must contain $0$. Direct sum: pairwise intersections not enough for $n\ge 3$.\
Basis $\neq$ ordered basis: order affects coordinates and matrices.\
Injective/surjective equivalence needs equal finite dimensions.\
Repeated eigenvalue does not settle diagonalizability; compute eigenspaces.\
Orthogonal $\neq$ orthonormal. Projection formulas require nonzero vector / ONB when using Fourier form.\
Over $\C$, keep conjugation in the correct slot.

\end{multicols}
\end{document}
